% Fibonacci-Pell nearest gaps: an all-exponent classification and
% quadratic-unit orbit rigidity
%
% Authors:
%   Dr. Denys Dutykh (Mathematics Department, Khalifa University of Science
%   and Technology, Abu Dhabi, UAE)
%   Prof. Laurent Vuillon (Univ. Savoie Mont Blanc, CNRS, LAMA, Chambery,
%   France)
%
\section{Local clocks, Haar recurrence, and their limits}
\label{sec:local-clocks}

The global classification uses the full norm--content equation. A tempting
alternative is to control only the common prime-power content forced by the
Fibonacci and Pell recurrences. We now show exactly how far that local idea
can go. Compatible $p$-adic clocks recur on positive-density subfamilies of
genuine nearest windows, but a rank-parity obstruction can also prevent a
split prime from occurring at all.

For the plus sign put

\begin{equation}
 D_q=F_q+F_{q+1}\sqrt2,
 \qquad
 m(q)=\left\lceil\log_\lambda D_q\right\rceil.
 \label{eq:local-m}
\end{equation}

Whenever $3\nmid q$ and

\begin{equation}
 D_q<\lambda^{m(q)}<\sqrt2D_q,
 \label{eq:local-window}
\end{equation}

the even power $\lambda^{2m(q)}$ lies in the strict nearest window for
$D_q^2$. Put $s=(-1)^{m(q)}$ and define

\begin{equation}
 G(q,m)=\gcd\left(
 R_q-s,\,
 2F_{q+1}P_m-F_qC_m
 \right).
 \label{eq:cross-gcd}
\end{equation}

This cross-gcd is defined on every row in \cref{eq:local-window}; the
norm--content equation is not assumed.

\begin{lemma}[Two exact phase factorizations]
\label{lem:cross-gcd-factorisations}
For $3\nmid q$,

\begin{equation}
\begin{aligned}
 m\equiv q\pmod2\quad&\Longrightarrow\quad
 G(q,m)=
 \gcd(F_{q+1},C_m)\gcd(F_{q+2},C_{m+1})\\[2mm]
 m\not\equiv q\pmod2\quad&\Longrightarrow\quad
 G(q,m)=
 \gcd(F_{q-1},C_{m-1})\gcd(F_{q+4},C_{m+2}).
\end{aligned}
\label{eq:cross-gcd-factorisations}
\end{equation}
\end{lemma}

\begin{proof}
Put $\varepsilon=(-1)^q$ and
$E=2F_{q+1}P_m-F_qC_m$. Cassini's identity gives

\begin{equation}
\begin{aligned}
 R_q-\varepsilon&=F_{q+1}F_{q+2},\\
 R_q+\varepsilon&=F_{q-1}F_{q+4}.
\end{aligned}
\label{eq:local-cassini-factors}
\end{equation}

Direct reduction gives

\begin{equation}
\begin{aligned}
 E&\equiv-F_qC_m&&\pmod {F_{q+1}},\\
 E&\equiv-F_qC_{m+1}&&\pmod {F_{q+2}},\\
 E&\equiv F_qC_{m-1}&&\pmod {F_{q-1}},\\
 3E&\equiv-F_qC_{m+2}&&\pmod {F_{q+4}}.
\end{aligned}
\label{eq:local-companion-reductions}
\end{equation}

The equal-phase Fibonacci factors are coprime. In the unequal phase,

$$
 \gcd(F_{q-1},F_{q+4})\in\{1,5\},
$$

but the exceptional prime $5$ divides no companion number $C_j$: the
recurrence modulo $5$ has period $12$ and nonzero residues
$1,1,3,2,2,1,4,4,2,3,3,4$. The only possible failure of the displayed
multipliers to be units is at $3$. Here

$$
 3\mid F_r\Longleftrightarrow4\mid r,
 \qquad
 3\mid C_j\Longleftrightarrow j\equiv2\pmod4,
$$

and every aligned Fibonacci--companion index pair has opposite parity.
Thus $3$ cannot occur in an aligned gcd. The product identity follows prime
by prime and valuation by valuation.
\end{proof}

The factorisation replaces one two-coordinate gcd by two aligned recurrence
clocks. We now determine exactly which primes can synchronize in one of
these clocks to arbitrary depth.

Let $p$ be an odd prime with $(2/p)=1$ that divides some $C_j$. Define

\begin{equation}
\begin{aligned}
 z_F(p)&=\min\{r>0:p\mid F_r\},
 &e_F(p)&=v_p(F_{z_F(p)}),\\
 z_C(p)&=\min\{r>0:p\mid C_r\},
 &e_C(p)&=v_p(C_{z_C(p)}),
\end{aligned}
\label{eq:local-ranks}
\end{equation}

and, for $a\geqslant1$,

\begin{equation}
 B_p(a)=z_F(p)p^{\max(0,a-e_F(p))},
 \qquad
 J_p(a)=z_C(p)p^{\max(0,a-e_C(p))}.
 \label{eq:lifted-ranks}
\end{equation}

The exact rank laws are

\begin{equation}
 p^a\mid F_r\Longleftrightarrow B_p(a)\mid r,
 \qquad
 p^a\mid C_j\Longleftrightarrow
 j\equiv J_p(a)\pmod {2J_p(a)}.
 \label{eq:rank-laws}
\end{equation}

The first is the standard Fibonacci rank and valuation law
\cite{Lengyel1995,Sanna2016}. For the second, choose a Hensel lift
$\sqrt2\in\Z_p$, put $u=1+\sqrt2$ and $b=-u^2$, and use

\begin{equation}
 2u^jC_j=(-1)^j(b^j+1).
 \label{eq:local-companion-unit}
\end{equation}

Companion zeros are the odd multiples of $z_C(p)$, and odd-prime lifting
gives their complete valuations. In local-field language,
$\Q(\sqrt2)\otimes\Q_p\simeq\Q_p\times\Q_p$ at such a split prime, and
\cref{eq:local-companion-unit} converts the clock into an order and
valuation statement for a unit of $\Z_p$; this is the special case of the
standard framework described by Serre \cite{Serre1979LocalFields}.

\begin{proposition}[Local-to-Archimedean transfer]
\label{prop:local-transfer}
Let $\mathcal P$ be a finite nonempty set of primes satisfying the preceding
hypotheses, prescribe $a_p\geqslant1$, and partition
$\mathcal P=\mathcal P_1\sqcup\mathcal P_2$. Choose
$\delta\in\{0,1\}$ and the offsets

\begin{equation}
\begin{array}{ccc}
 \toprule
 \delta&(r_1,t_1)&(r_2,t_2)\\
 \midrule
 0&(1,0)&(2,1)\\
 1&(-1,-1)&(4,2)\\
 \bottomrule
\end{array}
\label{eq:clock-offsets}
\end{equation}

Suppose that integers $q_0,m_0$ satisfy

\begin{equation}
 3\nmid q_0,
 \qquad
 q_0-m_0\equiv\delta\pmod2,
 \label{eq:clock-phase}
\end{equation}

and, for $i\in\{1,2\}$ and $p\in\mathcal P_i$,

\begin{equation}
 q_0+r_i\equiv0\pmod {B_p(a_p)},
 \qquad
 m_0+t_i\equiv J_p(a_p)\pmod {2J_p(a_p)}.
 \label{eq:clock-congruences}
\end{equation}

Put

\begin{equation}
 Q=\operatorname{lcm}\left(6,\{B_p(a_p):p\in\mathcal P\}\right),
 \qquad
 M=\operatorname{lcm}\left(\{2J_p(a_p):p\in\mathcal P\}\right).
 \label{eq:clock-moduli}
\end{equation}

Then infinitely many $q=q_0+Q\ell$ satisfy
\cref{eq:local-window}, $m(q)\equiv m_0\pmod M$, and

\begin{equation}
 \prod_{p\in\mathcal P}p^{a_p}\mid G(q,m(q)),
 \label{eq:clock-product}
\end{equation}

with every prime power carried by its prescribed factor in
\cref{eq:cross-gcd-factorisations}. Among the progression
$q=q_0+Q\ell$, this selected family has relative natural density

\begin{equation}
 \frac{\vartheta}{M},
 \qquad
 \vartheta=\log_\lambda\sqrt2.
 \label{eq:clock-relative-density}
\end{equation}

Its counting function up to $q\leqslant X$ is therefore
$\vartheta X/(MQ)+o(X)$.
\end{proposition}

\begin{proof}
The congruences and \cref{eq:rank-laws} place each prescribed prime power in
the required Fibonacci--companion pair of
\cref{eq:cross-gcd-factorisations}. They remain valid on the progression
$q=q_0+Q\ell$.

It remains to meet the Archimedean window with the required value of $m$.
Put

\begin{equation}
 \alpha=\frac{\log\varphi}{\log\lambda},
 \qquad
 \beta=\log_\lambda\left(\frac{1+\sqrt2\varphi}{\sqrt5}\right).
\end{equation}

Binet's formula gives

\begin{equation}
 \log_\lambda D_q=q\alpha+\beta+O(\varphi^{-2q}).
 \label{eq:local-rotation}
\end{equation}

The irrationality of $\alpha$ was proved in
\cref{thm:candidate-density}; hence $Q\alpha/M$ is irrational. Weyl
equidistribution on the circle $\R/M\Z$ shows that the limiting phase lies
in the cyclic interval $(m_0-\vartheta,m_0)$ with normalized Haar measure
$\vartheta/M$. Its boundary has Haar measure zero, and the Binet error tends
to zero. Membership in this interval is exactly
$m(q)\equiv m_0\pmod M$ together with \cref{eq:local-window}. This proves
the divisibility and both density assertions.
\end{proof}

\begin{theorem}[Exact clock criterion and simultaneous products]
\label{thm:local-clocks}
Let $p$ be an odd prime with $(2/p)=1$ that divides some $C_j$. Then the
following are equivalent:

\begin{enumerate}[label=\textup{(\roman*)}]
 \item $z_F(p)$ and $z_C(p)$ are not both even;
 \item for every $a\geqslant1$, infinitely many nearest-window rows have
 $p^a\mid G(q,m(q))$;
 \item at least one such row has $p\mid G(q,m(q))$.
\end{enumerate}

Every prime $p\equiv7\pmod8$ satisfies these conditions. Consequently, for
every finite set $\mathcal P$ of such primes and arbitrary exponents
$a_p\geqslant1$, infinitely many rows satisfy the one-bin divisibility

\begin{equation}
 \prod_{p\in\mathcal P}p^{a_p}
 \mid\gcd(F_{q+1},C_{m(q)})
 \mid G(q,m(q)).
 \label{eq:seven-class-product}
\end{equation}

In particular, for every $a,b\geqslant1$, infinitely many rows satisfy

\begin{equation}
 7^a17^b\mid G(q,m(q)).
 \label{eq:seven-seventeen}
\end{equation}

The rank-parity condition is sharp: $z_F(241)=120$ and $z_C(241)=20$, so

\begin{equation}
 241\nmid G(q,m)
 \label{eq:241-obstruction}
\end{equation}

for every plus nearest-window row.
\end{theorem}

\begin{proof}
For a single prime, use the first factor in the appropriate phase and write
$q=B_p(a)u-1$, with
$m\equiv J_p(a)\pmod {2J_p(a)}$. The phase condition becomes

\begin{equation}
 B_p(a)u\equiv J_p(a)+1\pmod2.
\end{equation}

If $B_p(a)$ is even, this is soluble exactly when $J_p(a)$ is odd; if
$B_p(a)$ is odd, the parity of $u$ is free. Multiplication by powers of the
odd prime does not change the rank parities, so the congruences are
consistent exactly when $z_F(p)$ and $z_C(p)$ are not both even. Choose the
residue of $u$ modulo $3$ so that $3\nmid B_p(a)u-1$. If
$3\mid B_p(a)$ this is automatic; otherwise exactly one residue is
forbidden. This choice is compatible with the required parity by the Chinese
remainder theorem modulo $6$.
\Cref{prop:local-transfer} then gives infinitely many rows at every depth.
Conversely, the Fibonacci and companion indices paired in
\cref{eq:cross-gcd-factorisations} have opposite parity. If both ranks are
even, no pair can contain $p$. This proves the equivalence.

Now let $p\equiv7\pmod8$ and choose $r^2=2\pmod p$. For
$u=1+r$, the order of $u$ has $2$-part at most $2$. If $h$ is its odd
part, then $u^h=\pm1$; since $1-r=-u^{-1}$,
$C_h\equiv0\pmod p$. Thus $z_C(p)$ is odd, and the criterion applies.
The assertion for $\mathcal P=\varnothing$ follows from
\cref{thm:candidate-density}, so suppose henceforth that
$\mathcal P\ne\varnothing$.
For finitely many such primes, the lifted companion ranks remain odd. Put

$$
 B=\operatorname{lcm}_{p\in\mathcal P}B_p(a_p),
 \qquad
 J=\operatorname{lcm}_{p\in\mathcal P}J_p(a_p).
$$

Then $m\equiv J\pmod {2J}$ satisfies every odd companion class, and
$q=Bu-1$ satisfies every first-bin Fibonacci class. Choose the parity of
$u$ to make $q\equiv m\pmod2$, and choose its residue modulo $3$ as above.
The choices are compatible modulo $6$, so \cref{prop:local-transfer} proves
\cref{eq:seven-class-product}. There are
infinitely many primes $7\pmod8$: if $p_1,\ldots,p_k$ were all of them,
an even multiple $A$ of their product would make $2A^2-1\equiv7\pmod8$;
every prime divisor $r$ has $(2/r)=1$, hence
$r\equiv1$ or $7\pmod8$, and at least one must be $7\pmod8$. It is new
because $2A^2-1\equiv-1\pmod {p_i}$ for every $i$.

For \cref{eq:seven-seventeen}, the exact lifted ranks are

\begin{equation}
\begin{aligned}
 B_7(a)&=8\mathbin{\cdot}7^{a-1},
 &J_7(a)&=3\mathbin{\cdot}7^{a-1},\\
 B_{17}(b)&=9\mathbin{\cdot}17^{b-1},
 &J_{17}(b)&=4\mathbin{\cdot}17^{b-1}.
\end{aligned}
\end{equation}

Assign $7$ and $17$ to the two equal-phase bins. Their $q$ moduli are
coprime and their $m$ moduli have gcd $2$ with odd prescribed residues, so
the systems are compatible. More explicitly,
$q\equiv-1\pmod {B_7(a)}$ makes $q$ odd, both companion congruences make
$m$ odd, and $q\equiv-2\pmod {B_{17}(b)}$ gives
$q\equiv1\pmod3$ because $9\mid B_{17}(b)$. Thus the phase and
$3\nmid q$ hypotheses of \cref{prop:local-transfer} hold. Finally, the
exact recurrence certificates at $241$ give the two even ranks displayed in
the theorem; the already proved criterion yields
\cref{eq:241-obstruction}.
\end{proof}

\begin{remark}[What the clocks establish]
\label{rem:local-scope}
The rows above satisfy genuine strict nearest windows, and the density is a
Haar-measure statement on a phase circle. They are not asserted to satisfy
\cref{eq:norm-content}. Thus local abundance produces neither a target nor a
Markoff counterexample. Its exact conclusion is methodological: uniform
boundedness, squarefreeness, finitely many compatible rank clocks, and the
assertion that every split prime occurs are all false as global substitutes
for the norm--content equation.
\end{remark}

This local--global contrast is the organising principle of the final
discussion, which situates the exact theorems relative to the parent Markoff
problem.
